\section{Introduction}
A service provider may observe many renewal histories while its terminal gateway accepts only a small menu of commands. The provider must choose that menu and allocate its promised service across histories. Participation limits expected service on each history, a separate ceiling limits each realized delivery, and installing a command can incur a fixed charge. Choosing a menu for prescribed branch targets is not the same optimization problem as choosing the menu and those targets jointly. This distinction matters most when the aggregate obligation is below the sum of the participation caps.

We develop an adaptive method for the joint finite-catalog problem. Its central ingredient is an exact price oracle on a box of target intervals. Unlike an unrestricted price problem, a box can force a branch to continue beyond the price-optimal command. Deleting all commands after the priced maximum can therefore destroy the correct upper bound. We show instead that the ordered menu has a two-sided representation: an increasing prefix accounts for upper-target crossings, a decreasing suffix accounts for lower-target crossings, and the common peak accounts for intervals containing it. Two coupled path recurrences optimize both sides at the original command budget, with every selected charge retained.

The resulting oracle yields globally valid bounds for adaptive target-space partitioning. The method returns a feasible contract and an exact interval containing the joint optimum even when stopped before reaching its requested accuracy. A zero-price witness makes the bound exact as a target box contracts. This supplies finite termination at any positive additive tolerance without completing a rectangular target grid. An independently implemented checker verifies all path inequalities, the entire target cover, and the original probability, service, promise, and realization constraints. It does not infer validity from an optimizer's completion flag.

Repeated histories can be combined when their caps and shortfall costs agree and they admit the same catalog prefix. Their realization ceilings need not be identical. We prove that this reduction preserves the exact joint optimum. The search dimension is consequently the number of distinct response types rather than the number of physical histories. This is an exact structural reduction, not approximate clustering. We also locate the elementary tractable regimes: one response type or a saturated obligation needs one conditional optimization; a fixed command budget admits polynomial catalog enumeration. We do not identify exponential dependence on the number of distinct types as a hardness lower bound.

The computational study evaluates joint, rather than solely conditional, design on a frozen synthetic holdout. Every completed and interrupted run is reported. The comparison includes exhaustive book search, a conventional subset-prefix baseline, independently formulated mixed-integer models, and the prior target grid. Larger repeated-type instances test the exact reduction explicitly; they are not evidence that arbitrary history populations possess a small number of types. The numerical section separates rational global intervals from floating-point solver bounds and from heuristic feasible values.

\paragraph{Relation to established methods.}\label{sec:r37-position}
Classical interval quantization and matrix searching provide the foundation for the saturated and conditional path algorithms \citep{Wu1991,NielsenNock2014,GronlundEtAl2017}. Pricing, constrained paths, and branch-and-bound are established methods \citep{HandlerZang1980,LawlerWood1966,Lemarechal2001}. Recent spatial branch-and-bound methods also exploit target intervals and locally tight relaxations for separable piecewise-linear optimization \citep{HubnerEtAl2025}. Our contribution is not the generic partitioning principle. It is the exact two-sided price decomposition for a shared, charged, cardinality-constrained menu under heterogeneous eligibility, and the root-feasible anchor condition that makes its box bound locally tight. The shared menu couples branches, so the joint objective is not a sum of independently specified piecewise-linear functions. The oracle handles that coupling through two path tables rather than introducing a separate binary decision for every branch--command pair.

The contractual and continuous-design results remain part of the analysis. Canonical allocation connects expected and pathwise institutions; the retained continuous quadratic and Monge results describe the design when command locations are not restricted to a catalog. Fixed-target compression, the full target net, distortion bounds, and charge-aware recovery remain available with their original assumptions. Table~\ref{tab:r46-contributions} distinguishes those antecedents from the present results. The main argument is now organized around joint design, while the companion retains the full preceding theorem chain and its computational evidence.
